关于神经网络最引人注目的事实之一是它们可以计算任意函数. 也就是说，假设有人交给你一个复杂的、弯弯曲曲的函数 $f(x)$：无论这个函数是什么，都保证存在一个神经网络，使得对于每一个可能的输入 $x$，网络都会输出值 $f(x)$（或某个接近的近似值），例如：即使函数有多个输入 $f = f(x_1, \ldots, x_m)$ 和多个输出，这一结果仍然成立. 例如，这里有一个计算具有 $m = 3$ 个输入和 $n = 2$ 个输出的函数的网络：这一结果告诉我们，神经网络具有一种\emph{通用性}（universality）\label{term:61}. 无论我们想要计算什么函数，我们都知道存在一个可以完成这项工作的神经网络.

更重要的是，即使我们将网络限制为在输入和输出神经元之间只有单层中间层——即所谓的单隐藏层，这一通用性定理仍然成立. 因此，即使是非常简单的网络结构也可以极其强大.

通用性定理对于使用神经网络的人来说是众所周知的. 但它为什么成立却并不那么广为人知. 大多数可获得的解释都相当技术化. 例如，证明这一结果的原始论文之一*\footnote{Approximation by superpositions of a sigmoidal function, by George Cybenko (1989). The result was very much in the air at the time, and several groups proved closely related results. Cybenko's paper contains a useful discussion of much of that work. Another important early paper is Multilayer feedforward networks are universal approximators, by Kurt Hornik, Maxwell Stinchcombe, and Halbert White (1989). This paper uses the Stone-Weierstrass theorem to arrive at similar results.} 使用了 Hahn-Banach 定理、Riesz 表示定理以及一些傅里叶分析来证明. 如果你是一个数学家，这个论证并不难理解，但对大多数人来说并不那么容易. 这很遗憾，因为通用性的根本原因是简单而优美的.

在本章中，我将给出一个简单且主要是可视化的通用性定理解释. 我们将一步一步地梳理其背后的思想. 你将理解为什么神经网络可以计算任意函数. 你将理解这一结果的一些局限性. 并且你将理解这一结果与深度神经网络的关系.

要跟上本章的内容，你不需要读过本书前面的章节. 相反，本章的结构使其作为一篇自成一体的文章而令人愉快. 只要你具备一点关于神经网络的基本知识，你就应该能够跟上解释. 不过，我会偶尔提供指向前面材料的链接，以帮助填补你知识中的任何空白.

通用性定理在计算机科学中是司空见惯的，以至于我们有时会忘记它们有多么令人惊叹. 但值得提醒我们自己：计算任意函数的能力确实是非凡的. 你能想象到的几乎任何过程都可以被视为函数计算. 考虑根据一小段音乐样本说出这段音乐的名称的问题. 这可以被视为计算一个函数. 或者考虑将一段中文文本翻译成英文的问题. 同样，这可以被视为计算一个函数*\footnote{实际上，是计算众多函数之一，因为对于给定的文本片段，通常有许多可接受的翻译.}. 或者考虑拿一个 mp4 电影文件，生成电影情节的描述，以及对表演质量的讨论的问题. 同样，这可以被视为一种函数计算*\footnote{关于翻译以及存在许多可能函数的评论同样适用.}. 通用性意味着，原则上，神经网络可以完成所有这些事情以及更多.

当然，仅仅因为我们知道存在一个可以（比如说）将中文文本翻译成英文的神经网络，并不意味着我们有好的技术来构造甚至识别这样一个网络. 这一局限性也适用于诸如布尔电路等模型的传统通用性定理. 但是，正如我们在本书前面所看到的，神经网络拥有强大的学习函数算法. 学习算法 + 通用性的组合是一种有吸引力的混合. 到目前为止，本书一直专注于学习算法. 在本章中，我们专注于通用性，以及它意味着什么.
\sectionTitle{两点说明}{两点说明}

在解释通用性定理为何成立之前，我想先提两点关于``神经网络可以计算任意函数''这一非正式说法的说明.

第一，这并不意味着网络能够被用来\emph{精确}计算任意函数. 更确切地说，我们可以得到一个\emph{近似}，并且这个近似可以任意好. 通过增加隐藏神经元的数量，我们可以改进近似效果. 例如，之前我用三个隐藏神经元展示了一个计算某个函数 $f(x)$ 的网络. 对于大多数函数来说，使用三个隐藏神经元只能得到低质量的近似. 通过增加隐藏神经元的数量（比如增加到五个），我们通常可以得到更好的近似：而进一步增加隐藏神经元的数量，我们还能做得更好.

为了更精确地表述这一点，假设给定一个函数 $f(x)$，我们希望以某个期望的精度 $\epsilon > 0$ 来计算它. 保证是：通过使用足够多的隐藏神经元，我们总能找到一个神经网络，其输出 $g(x)$ 对所有输入 $x$ 都满足 $|g(x) - f(x)| < \epsilon$. 换句话说，对于每一个可能的输入，近似都能达到期望的精度.

第二点说明是，能够以所述方式被近似的函数类，是\emph{连续}函数. 如果一个函数是不连续的，即发生突然的、急剧的跳跃，那么一般而言就不可能用神经网络来近似. 这并不奇怪，因为我们的神经网络计算的是其输入的连续函数. 然而，即使我们真正想要计算的函数是不连续的，也常常存在一个足够好的连续近似. 如果是这样，那么我们就可以使用神经网络. 在实践中，这通常不是一个重要的限制.

总结起来，通用性定理更精确的表述是：具有单个隐藏层的神经网络可以被用来以任意期望的精度近似任意连续函数. 在本章中，我们实际上将证明这个结果的一个稍弱版本，使用的是两个隐藏层而不是一个. 在习题中，我将简要概述如何稍作调整，就能使这一解释适用于给出一个只使用单个隐藏层的证明.
\sectionTitle{单输入单输出的通用性}{单输入单输出的通用性}

为了理解通用性定理为何成立，让我们先来理解如何构造一个仅有一个输入和一个输出的神经网络来近似一个函数：事实证明，这正是通用性问题的核心. 一旦我们理解了这个特殊情况，将其推广到具有多个输入和多个输出的函数实际上就相当容易了.

为了深入理解如何构造一个计算 $f$ 的网络，让我们从一个仅包含单个隐藏层、具有两个隐藏神经元，以及一个包含单个输出神经元的输出层的网络开始：为了感受网络中各个组件是如何工作的，让我们关注上面的隐藏神经元. 在下面的图中，点击权重 $w$，并将鼠标向右拖动一点以增大 $w$. 你可以立即看到上面隐藏神经元所计算的函数如何变化：正如我们在本书前面所学到的，隐藏神经元所计算的是 $\sigma(wx + b)$，其中 $\sigma(z) \equiv 1/(1+e^{-z})$ 是 S 型函数. 到目前为止，我们频繁地使用了这种代数形式. 但对于通用性的证明，我们将完全忽略代数，转而通过操作和观察图中所示的形状来获得更多的洞见. 这不仅会让我们更好地感受正在发生的事情，还会给我们一个适用于 S 型函数以外的激活函数的通用性证明*\footnote{严格来说，我所采用的视觉方法并不是传统意义上所认为的证明. 但我相信，视觉方法比传统证明更能让人洞见结果为何成立. 当然，这种洞见才是证明的真正目的. 偶尔，我所呈现的推理中会有一些小缺口：在这些地方，我给出了一个看似合理但不太严格的视觉论证. 如果这让你感到困扰，那么就把它当作一个填补缺失步骤的挑战. 但不要忘记真正的目的：理解通用性定理为何成立.}.

为了开始这个证明，请尝试点击上图中的偏置 $b$，并向右拖动以增大它. 你会看到，随着偏置增大，图形向左移动，但其形状不变.

接下来，点击并向左拖动以减小偏置. 你会看到，随着偏置减小，图形向右移动，但同样，其形状不变.

接下来，将权重减小到大约 $2$ 或 $3$. 你会看到，随着权重减小，曲线变宽. 你可能还需要改变偏置，以使曲线保持在视野内.

最后，将权重增大到超过 $w = 100$. 随着你这样做，曲线变得更陡峭，直到最终开始看起来像一个阶梯函数. 尝试调整偏置，使阶梯出现在 $x = 0.3$ 附近. 下面的短片展示了你的结果应该是什么样子. 点击播放按钮来播放（或重播）视频：

\begin{movieNote}{create\_\allowbreak{}step\_\allowbreak{}function}
\end{movieNote}

\noindent 我们可以通过将权重增大到输出在非常好的近似下确实是一个阶梯函数，来大大简化我们的分析. 下面我绘制了当权重为 $w = 999$ 时上面隐藏神经元的输出. 请注意，这个图是静态的，你无法改变权重等参数.

实际上，使用阶梯函数比使用一般的 S 型函数要容易得多. 原因在于，在输出层中，我们将所有隐藏神经元的贡献相加. 分析一堆阶梯函数的和很容易，但推理当你把一堆 S 型曲线相加时会发生什么则要困难得多. 因此，假设我们的隐藏神经元输出阶梯函数会使事情容易得多. 更具体地说，我们通过将权重 $w$ 固定为某个非常大的值，然后通过修改偏置来设置阶梯的位置来做到这一点. 当然，将输出视为阶梯函数是一种近似，但这是一个非常好的近似，目前我们将把它视为精确的. 稍后我会回来讨论偏离这种近似的影响.

阶梯出现在 $x$ 的什么值处？换句话说，阶梯的位置如何依赖于权重和偏置？

为了回答这个问题，请尝试修改上图中的权重和偏置（你可能需要向上滚动一点）. 你能弄清楚阶梯的位置如何依赖于 $w$ 和 $b$ 吗？稍加努力，你应该能够说服自己，阶梯的位置与 $b$ \emph{成正比}，与 $w$ \emph{成反比}.

事实上，阶梯位于位置 $s = -b/w$，你可以通过修改下图中的权重和偏置来看到这一点：仅使用单个参数 $s$ 来描述隐藏神经元将大大简化我们的生活，$s$ 就是阶梯位置，$s = -b/w$. 尝试修改下图中的 $s$，以习惯新的参数化方式：如上所述，我们隐含地将输入的权重 $w$ 设置为某个很大的值——大到足以使阶梯函数成为一个非常好的近似. 我们可以很容易地将以这种方式参数化的神经元转换回传统模型，只需选择偏置 $b = -w s$.

到目前为止，我们一直关注的是仅上面隐藏神经元的输出. 让我们来看看整个网络的行为. 特别地，我们将假设隐藏神经元计算由阶梯点 $s_1$（上面的神经元）和 $s_2$（下面的神经元）参数化的阶梯函数. 它们将分别具有输出权重 $w_1$ 和 $w_2$. 这是网络：右边绘制的是来自隐藏层的\emph{加权输出} $w_1 a_1 + w_2 a_2$. 这里，$a_1$ 和 $a_2$ 分别是上面和下面隐藏神经元的输出*\footnote{顺便注意，整个网络的输出是 $\sigma(w_1 a_1+w_2 a_2 + b)$，其中 $b$ 是输出神经元的偏置. 显然，这与我们这里绘制的来自隐藏层的加权输出不同. 我们现在将专注于来自隐藏层的加权输出，稍后才会考虑它与整个网络输出的关系.}. 这些输出用 $a$ 表示，因为它们通常被称为神经元的\emph{激活值}.

尝试增大和减小上面隐藏神经元的阶梯点 $s_1$. 感受一下这如何改变来自隐藏层的加权输出. 特别值得理解的是当 $s_1$ 越过 $s_2$ 时会发生什么. 你会看到，当这种情况发生时，图形会改变形状，因为我们已经从上面隐藏神经元最先被激活的情况转变为下面隐藏神经元最先被激活的情况.

同样，尝试操作下面隐藏神经元的阶梯点 $s_2$，并感受一下这如何改变来自隐藏神经元的组合输出.

尝试增大和减小每个输出权重. 注意这如何重新缩放来自相应隐藏神经元的贡献. 当其中一个权重为零时会发生什么？

最后，尝试将 $w_1$ 设为 $0.8$，将 $w_2$ 设为 $-0.8$. 你会得到一个``凸起''函数，它从点 $s_1$ 开始，到点 $s_2$ 结束，高度为 $0.8$. 例如，加权输出可能看起来像这样：当然，我们可以重新缩放凸起使其具有任意高度. 让我们使用单个参数 $h$ 来表示高度. 为了减少杂乱，我还将移除``$s_1 = \ldots$''和``$w_1 = \ldots$''的标注.

尝试上下改变 $h$ 的值，看看凸起的高度如何变化. 尝试将高度改为负值，并观察会发生什么. 并尝试改变阶梯点，看看这如何改变凸起的形状.

顺便你会注意到，我们使用神经元的方式不仅可以从图形角度来思考，还可以从更传统的编程角度来思考，作为一种 \texttt{if-then-else} 语句，例如：

\begin{lstlisting}
    if input >= step point:
        add 1 to the weighted output
    else:
        add 0 to the weighted output

\end{lstlisting}

在大多数情况下，我将坚持图形视角. 但在接下来的内容中，你有时可能会发现切换视角、从 \texttt{if-then-else} 的角度来思考事情会有所帮助.

我们可以使用制造凸起的技巧来得到两个凸起，方法是将两对隐藏神经元粘合到同一个网络中：我在这里省略了权重，只写下每对隐藏神经元的 $h$ 值. 尝试增大和减小两个 $h$ 值，并观察它如何改变图形. 通过改变阶梯点来移动凸起.

更一般地，我们可以使用这个想法来获得任意数量的任意高度的峰. 特别地，我们可以将区间 $[0, 1]$ 分成大量 $N$ 个子区间，并使用 $N$ 对隐藏神经元来设置任意所需高度的峰. 让我们看看这对于 $N = 5$ 是如何工作的. 那是相当多的神经元，所以我要把东西压缩一下. 对图的复杂性表示歉意：我可以通过进一步抽象来隐藏复杂性，但我认为为了更具体地感受这些网络如何工作，忍受一点复杂性是值得的.

你可以看到有五对隐藏神经元. 各对神经元的阶梯点分别是 $0, 1/5$，然后是 $1/5, 2/5$，依此类推，一直到 $4/5, 5/5$. 这些值是固定的——它们使我们得到图上五个均匀间隔的凸起.

每对神经元都有一个与之关联的 $h$ 值. 记住，从神经元输出的连接具有权重 $h$ 和 $-h$（未标出）. 点击其中一个 $h$ 值，并将鼠标向右或向左拖动以改变该值. 当你这样做时，观察函数的变化. 通过改变输出权重，我们实际上是在\emph{设计}函数！

反之，尝试点击图形，并向上或向下拖动以改变任何凸起函数的高度. 当你改变高度时，你可以看到 $h$ 值的相应变化. 而且，虽然没有显示出来，相应的输出权重也会发生变化，即 $+h$ 和 $-h$.

换句话说，我们可以直接操作右侧图形中出现的函数，并在左侧的 $h$ 值中看到其反映. 一个有趣的事情是按住鼠标按钮，将鼠标从图形的一侧拖到另一侧. 当你这样做时，你画出一个函数，并可以观察神经网络中的参数随之调整.

挑战时间.

让我们回想一下我在本章开头绘制的函数：我当时没有说，但我绘制的实际上是函数

\begin{equation}
f(x) = 0.2+0.4 x^2+0.3x \sin(15 x) + 0.05 \cos(50 x),
\tag{113}
\end{equation}

\noindent 在 $x$ 从 $0$ 到 $1$ 上绘制，$y$ 轴取值从 $0$ 到 $1$.

这显然不是一个平凡的函数.

你将弄清楚如何使用神经网络来计算它.

在我们上面的网络中，我们一直在分析从隐藏神经元输出的加权组合 $\sum_j w_j a_j$. 我们现在知道如何对这个量获得很多控制. 但是，正如我之前所指出的，这个量不是网络的输出. 网络的输出是 $\sigma(\sum_j w_j a_j + b)$，其中 $b$ 是输出神经元的偏置. 有没有什么方法我们可以实现对网络实际输出的控制？

解决方案是设计一个神经网络，其隐藏层的加权输出由 $\sigma^{-1} \circ f(x)$ 给出，其中 $\sigma^{-1}$ 就是 $\sigma$ 函数的反函数. 也就是说，我们希望来自隐藏层的加权输出为：如果我们能做到这一点，那么整个网络的输出将很好地近似 $f(x)$*\footnote{注意，我已将输出神经元的偏置设为 $0$.}.

那么，你的挑战是设计一个神经网络来近似上面显示的目标函数\label{term:62}. 为了尽可能多地学习，我希望你解决这个问题两次. 第一次，请点击图形，直接调整不同凸起函数的高度. 你应该会发现相当容易获得与目标函数的良好匹配. 你的表现如何由目标函数与网络实际计算的函数之间的\emph{平均偏差}来衡量. 你的挑战是尽可能\emph{低}地驱动平均偏差. 当你将平均偏差降到 $0.40$ 或以下时，你就完成了挑战.

一旦你做到了这一点，点击``Reset''来随机重新初始化凸起. 第二次解决这个问题时，请抵制点击图形的冲动. 相反，修改左侧的 $h$ 值，并再次尝试将平均偏差降到 $0.40$ 或以下.

你现在已经弄清楚了网络近似计算函数 $f(x)$ 所需的所有要素！这只是一个粗略的近似，但我们可以很容易地做得更好，只需增加隐藏神经元的对数，允许更多的凸起.

特别地，很容易将我们找到的所有数据转换回用于神经网络的标准参数化方式. 让我快速回顾一下这是如何工作的.

第一层权重都具有某个很大的常数值，比如 $w = 1000$.
隐藏神经元上的偏置就是 $b = -w s$. 因此，例如，对于第二个隐藏神经元 $s = 0.2$ 变为 $b = -1000 \times 0.2 = -200$.

最后一层的权重由 $h$ 值确定. 因此，例如，你上面为第一个 $h$ 选择的值 $h = $，意味着来自顶部两个隐藏神经元的输出权重分别为 和. 以此类推，对于整个输出权重层.

最后，输出神经元上的偏置为 $0$.

这就是全部内容：我们现在有了一个神经网络的完整描述，它能够相当好地计算我们最初的目标函数. 并且我们理解了如何通过增加隐藏神经元的数量来提高近似的质量.

更重要的是，我们最初的目标函数 $f(x) = 0.2+0.4 x^2+0.3 \sin(15 x) + 0.05 \cos(50 x)$ 并没有什么特别之处. 我们本可以将这个过程用于任何从 $[0, 1]$ 到 $[0, 1]$ 的连续函数. 本质上，我们是在用我们的单层神经网络为函数构建一个查找表. 并且我们将能够在这一想法的基础上，给出通用性的一个一般性证明.
\sectionTitle{多个输入变量}{多个输入变量}

让我们把结果推广到多个输入变量的情形. 这听起来很复杂，但我们所需要的全部思想都可以在两个输入的情形下理解. 所以让我们先处理两个输入的情形.

我们先来考虑一个神经元有两个输入时会发生什么：这里，我们有输入 $x$ 和 $y$，对应的权重为 $w_1$ 和 $w_2$，以及神经元上的偏置 $b$. 让我们把权重 $w_2$ 设为 $0$，然后摆弄第一个权重 $w_1$ 和偏置 $b$，看看它们如何影响神经元的输出：如你所见，当 $w_2 = 0$ 时，输入 $y$ 对神经元的输出没有任何影响. 就好像 $x$ 是唯一的输入.

既然如此，你认为当我们把权重 $w_1$ 增大到 $w_1 = 100$，而 $w_2$ 保持为 $0$ 时会发生什么？如果你没有立刻看出答案，就思考这个问题一会儿，看看你能否弄清楚会发生什么. 然后试一试，看看你是否正确. 我在下面的影片中展示了会发生什么：

\begin{movieNote}{step\_\allowbreak{}3d}
\end{movieNote}

\noindent 正如我们之前讨论的那样，随着输入权重变大，输出趋近于一个阶梯函数. 不同之处在于，现在的阶梯函数是三维的. 同样和之前一样，我们可以通过修改偏置来移动阶梯点的位置. 阶梯点的实际位置是 $s_x \equiv -b / w_1$.

让我们用阶梯的位置作为参数重做上面的过程：这里，我们假设 $x$ 输入上的权重有一个很大的值——我用了 $w_1 = 1000$——并且权重 $w_2 = 0$. 神经元上的数字是阶梯点，数字上方的小 $x$ 提醒我们阶梯是在 $x$ 方向上的. 当然，也可以通过让 $y$ 输入上的权重非常大（比如 $w_2 = 1000$），并让 $x$ 上的权重等于 $0$，即 $w_1 = 0$，来得到一个 $y$ 方向上的阶梯函数：神经元上的数字同样是阶梯点，在这种情况下数字上方的小 $y$ 提醒我们阶梯是在 $y$ 方向上的. 我本可以明确标出 $x$ 和 $y$ 输入上的权重，但决定不这样做，因为那会让图变得相当杂乱. 但请记住，小 $y$ 标记隐含地告诉我们 $y$ 权重很大，而 $x$ 权重为 $0$.

我们可以用刚刚构造的阶梯函数来计算一个三维的凸起函数. 为此，我们使用两个神经元，每个计算一个 $x$ 方向上的阶梯函数. 然后我们分别用权重 $h$ 和 $-h$ 组合这些阶梯函数，其中 $h$ 是凸起所期望的高度. 这一切都在下面的图中展示：试着改变高度 $h$ 的值. 观察它如何与网络中的权重相关联. 并看看它如何改变右边凸起函数的高度.

另外，试着改变与顶部隐藏神经元相关联的阶梯点 $0.30$. 观察它如何改变凸起的形状. 当你把它移过与底部隐藏神经元相关联的阶梯点 $0.70$ 时会发生什么？

我们已经弄清楚了如何在 $x$ 方向上制作一个凸起函数. 当然，我们可以很容易地在 $y$ 方向上制作一个凸起函数，方法是在 $y$ 方向上使用两个阶梯函数. 回想一下，我们通过让 $y$ 输入上的权重很大，并让 $x$ 输入上的权重为 $0$ 来做到这一点. 结果如下：这看起来与之前的网络几乎完全相同！唯一明确显示为变化的是，我们的隐藏神经元上现在有小 $y$ 标记. 这提醒我们它们产生的是 $y$ 阶梯函数，而不是 $x$ 阶梯函数，因此 $y$ 输入上的权重非常大，而 $x$ 输入上为零，而不是反过来. 和之前一样，我决定不明确展示这一点，以避免杂乱.

让我们考虑当把两个凸起函数相加时会发生什么，一个在 $x$ 方向上，另一个在 $y$ 方向上，高度都为 $h$：为了简化图，我省略了权重为零的连接. 目前，我保留了隐藏神经元上的小 $x$ 和 $y$ 标记，以提醒你凸起函数是在哪些方向上计算的. 稍后我们会连这些标记也去掉，因为它们已经由输入变量隐含了.

试着改变参数 $h$. 如你所见，这会导致输出权重发生变化，同时也会改变 $x$ 和 $y$ 凸起函数的高度.

我们构建的东西看起来有点像 \emph{塔} 函数：如果我们能构建这样的塔函数，那么我们就可以用它们来近似任意函数，只需把许多不同高度、不同位置的塔加起来：当然，我们还没有弄清楚如何构建塔函数. 我们构造出来的东西看起来像一个中心塔，高度为 $2h$，周围有一个高度为 $h$ 的平台.

但我们可以制作一个塔函数. 记住，之前我们看到神经元可以用来实现一种 \texttt{if-then-else} 语句：

\begin{lstlisting}
    if input >= threshold:
        output 1
    else:
        output 0

\end{lstlisting}

那是一个只有一个输入的神经元. 我们想要的是把类似的想法应用到隐藏神经元的组合输出上：

\begin{lstlisting}
    if combined output from hidden neurons >= threshold:
        output 1
    else:
        output 0

\end{lstlisting}

如果我们适当地选择 \texttt{threshold}——比如取值为 $3h/2$，它夹在平台的高度和中心塔的高度之间——我们就可以把平台压到零，只留下塔立着.

你能看出如何做到这一点吗？试着在下面的网络中做实验来弄清楚. 注意我们现在绘制的是整个网络的输出，而不仅仅是隐藏层的加权输出. 这意味着我们在隐藏层的加权输出上加上一个偏置项，并应用 sigma 函数. 你能找到产生塔的 $h$ 和 $b$ 的值吗？这有点棘手，所以如果你思考了一会儿仍然卡住，这里有两个提示：(1) 要让输出神经元表现出正确类型的 \texttt{if-then-else} 行为，我们需要输入权重（全部为 $h$ 或 $-h$）很大；(2) $b$ 的值决定了 \texttt{if-then-else} 阈值的尺度.

用我们初始的参数，输出看起来像是之前图的扁平版本，有它的塔和平台. 为了得到期望的行为，我们增大参数 $h$ 直到它变得很大. 这就给出了 \texttt{if-then-else} 阈值行为. 其次，为了把阈值调对，我们选择 $b \approx -3h/2$. 试试看，看看效果如何！

当我们使用 $h = 10$ 时，它看起来是这样的：

\begin{movieNote}{tower\_\allowbreak{}construction}
\end{movieNote}

\noindent 即使对于这个相对适中的 $h$ 值，我们也得到了一个相当不错的塔函数. 而且，当然，我们可以通过进一步增大 $h$，并保持偏置为 $b = -3h/2$，让它变得要多好有多好.

让我们试着把两个这样的网络粘在一起，以计算两个不同的塔函数. 为了清楚表明两个子网络各自的作用，我把它们放在下面单独的框中：每个框使用上述技术计算一个塔函数. 右边的图显示了 \emph{第二个} 隐藏层的加权输出，也就是说，它是塔函数的加权组合.

特别地，你可以看到通过修改最后一层的权重，你可以改变输出塔的高度.

同样的想法可以用来计算任意多个塔. 我们也可以让它们任意细，任意高. 因此，我们可以确保第二个隐藏层的加权输出近似任意期望的二元函数：特别地，通过让第二个隐藏层的加权输出很好地近似 $\sigma^{-1} \circ f$，我们确保网络的输出将很好地近似任意期望的函数 $f$.

多于两个变量的函数呢？

让我们试试三个变量 $x_1, x_2, x_3$. 下面的网络可以用来计算一个四维的塔函数：这里，$x_1, x_2, x_3$ 表示网络的输入. $s_1, t_1$ 等等是神经元的阶梯点——也就是说，第一层中的所有权重都很大，偏置被设置为给出阶梯点 $s_1, t_1, s_2, \ldots$. 第二层中的权重交替为 $+h, -h$，其中 $h$ 是某个非常大的数. 输出偏置为 $-5h/2$.

这个网络计算一个函数，当满足三个条件时它为 $1$：$x_1$ 在 $s_1$ 和 $t_1$ 之间；$x_2$ 在 $s_2$ 和 $t_2$ 之间；$x_3$ 在 $s_3$ 和 $t_3$ 之间. 网络在其他任何地方都为 $0$. 也就是说，它是一种塔，在输入空间的一小块区域中为 $1$，在其他任何地方为 $0$.

通过把许多这样的网络粘在一起，我们可以得到任意多个塔，从而近似任意三元函数. 完全相同的想法在 $m$ 维中也适用. 唯一需要的变化是把输出偏置设为 $(-m+1/2)h$，以获得正确类型的夹逼行为来压平平台.

好了，所以我们现在知道如何使用神经网络来近似多变量的实值函数. 那么向量值函数 $f(x_1, \ldots, x_m) \in R^n$ 呢？当然，这样的函数可以看作只是 $n$ 个独立的实值函数，$f^1(x_1, \ldots, x_m), f^2(x_1, \ldots, x_m)$，等等. 所以我们创建一个近似 $f^1$ 的网络，另一个用于 $f^2$ 的网络，等等. 然后我们简单地把所有网络粘在一起. 所以这也很容易处理.

\begin{exerciseBox}{习题}

\begin{itemize}

\item 我们已经看到了如何使用带有两个隐藏层的网络来近似任意函数. 你能找到一个证明，表明只用单个隐藏层也是可能的吗？提示：试着在两个输入变量的情形下工作，并证明：(a) 不仅可以得到 $x$ 或 $y$ 方向上的阶梯函数，还可以得到任意方向上的阶梯函数；(b) 通过把 (a) 中的许多构造加起来，可以近似一个形状为圆形而非矩形的塔函数；(c) 使用这些圆形塔，可以近似任意函数. 做 (c) 部分时，使用本章稍后的一些想法可能会有所帮助.

\end{itemize}

\end{exerciseBox}
\sectionTitle{超越 S 型神经元}{超越 S 型神经元}

我们已经证明，由 S 型神经元构成的网络可以计算任意函数. 回忆一下，在 S 型神经元中，输入 $x_1, x_2, \ldots$ 产生输出 $\sigma(\sum_j w_j x_j + b)$，其中 $w_j$ 是权重，$b$ 是偏置，$\sigma$ 是 S 型函数：如果我们考虑一种不同类型的神经元，使用其他某个激活函数 $s(z)$，会怎样呢：也就是说，我们将假设如果我们的神经元有输入 $x_1, x_2, \ldots$、权重 $w_1, w_2, \ldots$ 和偏置 $b$，那么输出就是 $s(\sum_j w_j x_j + b)$.

我们可以用这个激活函数得到一个阶梯函数，就像我们对 S 型函数所做的那样. 试着在下面把权重调大，比如调到 $w = 100$：就像 S 型函数一样，这会使激活函数收缩，最终它成为阶梯函数的一个非常好的近似. 试着改变偏置，你会看到我们可以把阶梯的位置设置在我们选择的任何地方. 因此我们可以使用和之前完全相同的技巧来计算任意想要的函数.

为了让这一切成立，$s(z)$ 需要满足什么性质呢？我们确实需要假设当 $z \rightarrow -\infty$ 和 $z \rightarrow \infty$ 时 $s(z)$ 有良好定义. 这两个极限就是我们的阶梯函数所取的两个值. 我们还需要假设这两个极限彼此不同. 如果它们相同，就不会有阶梯，而只是一条平坦的曲线！但只要激活函数 $s(z)$ 满足这些性质，基于这种激活函数的神经元对于计算就是通用的.

\begin{exerciseBox}{习题}

\begin{itemize}

\item 在本书前面我们遇到了另一种神经元，称为修正线性单元. 解释为什么这种神经元不满足刚刚给出的通用性条件. 找出一个通用性证明，表明修正线性单元对于计算是通用的.

\item 假设我们考虑线性神经元，即激活函数为 $s(z) = z$ 的神经元. 解释为什么线性神经元不满足刚刚给出的通用性条件. 证明这种神经元不能用于进行通用计算.

\end{itemize}

\end{exerciseBox}
\sectionTitle{修补阶梯函数}{修补阶梯函数}

到目前为止，我们一直假设我们的神经元能够精确地产生阶梯函数. 这是一个相当好的近似，但它只是一个近似. 事实上，会存在一个狭窄的失败窗口，如下图所示，在这个窗口中，函数的行为与阶梯函数非常不同：在这些失败窗口中，我给出的通用性解释将会失效.

现在，这还不是一个糟糕的失败. 通过使输入到神经元的权重足够大，我们可以使这些失败窗口变得任意小. 当然，我们可以使窗口比我上面展示的窄得多——事实上，比我们的眼睛能看到的还要窄. 所以也许我们不必太担心这个问题.

尽管如此，如果能有一些方法来解决这个问题，那将是很好的.

事实上，这个问题很容易解决. 让我们看看针对只有一个输入和一个输出的神经网络计算函数的修复方法. 同样的想法也适用于有更多输入和输出的情况.

特别地，假设我们希望我们的网络计算某个函数 $f$. 和之前一样，我们通过尝试设计我们的网络，使得来自隐藏层神经元的加权输出为 $\sigma^{-1} \circ f(x)$ 来实现这一点：如果我们使用之前描述的技术来做这件事，我们会使用隐藏神经元来产生一系列隆起函数：同样，我夸大了失败窗口的大小，以便更容易看到它们. 应该很清楚，如果我们将所有这些隆起函数加起来，我们最终会得到一个对 $\sigma^{-1} \circ f(x)$ 的合理近似，除了在失败窗口内.

假设我们不是使用刚刚描述的近似，而是使用一组隐藏神经元来计算我们原始目标函数的\emph{一半}的近似，即 $\sigma^{-1} \circ f(x) / 2$. 当然，这看起来就像是最后一个图形的缩小版本：并且假设我们使用另一组隐藏神经元来计算 $\sigma^{-1} \circ f(x)/ 2$ 的近似，但隆起的基底\emph{偏移}了半个隆起的宽度：现在我们有了两个不同的对 $\sigma^{-1} \circ f(x) / 2$ 的近似. 如果我们将这两个近似加起来，我们将得到对 $\sigma^{-1} \circ f(x)$ 的整体近似. 这个整体近似仍然会在小窗口中有失败. 但问题将比以前小得多. 原因是一个近似的失败窗口中的点不会在另一个近似的失败窗口中. 因此，在这些窗口中，近似将大约好 $2$ 倍.

我们可以通过将大量 $M$ 个重叠的近似加到函数 $\sigma^{-1} \circ f(x) / M$ 上来做得更好. 只要失败窗口足够窄，一个点将只会出现在一个失败窗口中. 并且只要我们使用足够大的数量 $M$ 个重叠近似，结果将是一个极好的整体近似.
\sectionTitle{结语}{结语}

我们讨论的通用性解释当然不是如何用神经网络进行计算的实用处方！在这一点上，它很像 \texttt{NAND} 门等的通用性证明. 出于这个原因，我主要关注的是让构造过程清晰易懂，而不是优化构造的细节. 不过，你可能会发现，看看自己能否改进这个构造，是一个有趣且有启发性的练习.

尽管这个结果在构造网络时没有直接用处，但它很重要，因为它排除了某个特定函数是否可以用神经网络计算的问题. 这个问题的答案总是``是''. 所以正确的问题不是某个特定函数是否可计算，而是计算该函数的\emph{好}方法是什么.

我们发展的通用性构造只用了两个隐藏层来计算任意函数. 此外，正如我们讨论过的，仅用一个隐藏层也有可能得到相同的结果. 既然如此，你可能会想，为什么我们还会对深度网络感兴趣，也就是具有许多隐藏层的网络. 难道我们不能简单地把那些网络替换成浅层的、单个隐藏层的网络吗？\footnote{\textbf{章节致谢：}感谢 Jen Dodd 和 Chris Olah 关于神经网络通用性的多次讨论. 我尤其要感谢 Chris 建议使用查找表来证明通用性. 本章的交互式可视化形式受到了 Mike Bostock、Amit Patel、Bret Victor 和 Steven Wittens 等人工作的启发.}虽然原则上这是可能的，但使用深度网络有很好的实际理由. 正如第 1 章所论证的，深度网络具有层次结构，这使得它们特别适合学习那些似乎在解决现实世界问题时有用的知识层次. 更具体地说，当处理图像识别等问题时，使用一个不仅理解单个像素，而且理解越来越复杂概念的系统会有所帮助：从边缘到简单的几何形状，一直到复杂的多物体场景. 在后面的章节中，我们将看到证据表明，深度网络在学习这类知识层次方面比浅层网络做得更好. 总结一下：通用性告诉我们神经网络可以计算任意函数；而经验证据表明，深度网络是最适合学习那些在解决许多现实世界问题中有用的函数的网络.
